\documentclass[10pt,a4paper]{amsart}
\usepackage[margin=19mm]{geometry}
\usepackage{amsmath,amssymb,mathtools}
\usepackage[scr=rsfs,cal=euler]{mathalfa}
\usepackage{xfrac}
\usepackage[unicode,hidelinks,bookmarks=false]{hyperref}
\newcommand{\ie}{i.e.\ }
\newcommand{\cf}{cf.\ }
\newcommand{\resp}{respectively\ }
\newcommand{\Subsection}{\subsection}
\providecommand{\nobreakdash}{\nobreak\hbox{-}}
\setlength{\emergencystretch}{2em}
\multlinegap0pt
\theoremstyle{plain}
\newtheorem{theorem}{Theorem}
\newtheorem{lemma}[theorem]{Lemma}
\newtheorem{proposition}[theorem]{Proposition}
\newtheorem{corollary}[theorem]{Corollary}
\theoremstyle{definition}
\newtheorem{definition}[theorem]{Definition}
\newcommand{\Bc}{\mathcal{B}}
\DeclareMathOperator{\sgn}{sgn}
\usepackage{cleveref}
\title[Boolean Rank via Monomial Ideals]{Boolean Rank via Monomial Ideals: monomials in the determinantal ideal correspond to isolated sets (printed as Theorem 4.13)}
\author{Juliann Geraci, Alexander B. Kunin, Alexandra Seceleanu}
\date{Source: arXiv:2509.09570v2 (CC BY 4.0)}
\begin{document}
\maketitle

\noindent\emph{Source and licence.} Boolean Rank via Monomial Ideals. Version arXiv:2509.09570v2 (CC BY 4.0), distributed under the Creative Commons Attribution 4.0 International licence, http://creativecommons.org/licenses/by/4.0/, as recorded in the arXiv licence metadata for this exact version and shown on the saved abstract page https://arxiv.org/abs/2509.09570v2. Full licence notice: licenses/arxiv-2509.09570-CC-BY-4.0.txt.\par\smallskip

\noindent\textit{Editorial scope.} Counted unit: the correspondence between monomials lying in a determinantal ideal of a Boolean matrix and isolated sets of entries, printed in the article as Theorem 4.13 (source0.tex lines 653--655, one \texttt{theorem} environment with source label \texttt{thm: monomials=isolated}), delivered together with the article's complete original proof of it (source0.tex lines 656--704, one \texttt{proof} environment of 49 lines immediately adjacent to the statement, with no gap and no deferral). The proof is an induction on the degree of the monomial generator whose induction step is carried out by expanding the relevant minor along its determinant, showing that every non-identity permutation term vanishes, and concluding that the smaller minor still lies in the smaller determinantal ideal before applying the induction hypothesis. No mathematics is written, repaired or re-derived: statement and proof are the article's own text line for line, in the article's own notation, and no line of either is omitted, substituted or re-flowed.

Required context retained uncounted: the article's Lemma 4.12 with its label (source0.tex lines 616--618) and the article's complete original proof of that lemma (source0.tex lines 619--649). The counted proof uses that lemma for its base case, so both the statement and the proof of the lemma are delivered; the two displays the induction step refers to are labelled inside the counted proof itself.

References. No citation command occurs in any retained line, so this package carries no bibliography block and no bibliography entry.

Editorial formatting only, in addition: an AMS \texttt{amsart} document, the same class the article itself uses, that restates the macro definitions the retained lines use (the article's own $\Bc$ and its $\sgn$ operator), loads the cleveref package for the article's own \texttt{\textbackslash Cref} cross-reference command, and restates the article's own theorem-environment declarations, keeping the article's convention that the lemma, proposition, corollary and definition environments share the theorem counter. Consequently the retained environments print here in the order in which they occur, so the retained lemma prints as Lemma 1 and the counted theorem as Theorem 2, whereas the article prints them as Lemma 4.12 and Theorem 4.13; because those two environments share one counter, the article's own \texttt{\textbackslash Cref} command quotes that counter's name and therefore reads Theorem 1 here and Theorem 4.12 in the article, which is exactly what the article does; the article's own per-section numbering is not reproduced, the equations of the retained spans are numbered anew in order of appearance, and every cross-reference inside the retained text resolves to this wrapper's numbering. That is the only change to the numbering. No figure, no image-inclusion line and no drawing code occur in any retained line, so no artwork is delivered or needed, although the article contains figures elsewhere.

The article's own title line reads Boolean Matrix Rank via Monomial Ideals, while the arXiv record for this version and the licence metadata give the title as Boolean Rank via Monomial Ideals, which is the form used here and throughout this package. The article has no separate class or style file; the package ships the article's concatenated source verbatim as \texttt{sources/184918/source0.tex}, and every delivered excerpt is an exact line range of that file. Not retained: the article's definitions, its other results, its examples, its computations and its bibliography; none of that material is used as a step of the delivered proof, which is complete as the author wrote it. No mathematical statement, hypothesis, estimate or conclusion has been rewritten, repaired or added, and printed slips, if any, are preserved literally.
\begin{lemma}\label{lem: 2x2 minors}
    Let $A\in \Bc^{m,n}$. Entries $a_{ij}$ and $a_{k\ell}$ are an isolated pair if and only if $x_{ij}x_{k\ell}\in I_2(A)$. 
\end{lemma}

\begin{proof}
Suppose $a_{ij}$ and $a_{k\ell}$ are an isolated pair, % we know that $a_{ij}a_{k\ell} = 1$ and $a_{i\ell}a_{kj}=0$. So this means that $a_{ij}=a_{k\ell}=1$ and either $a_{i\ell}=0$ or $a_{kj}=0$ (possibly both). Then in $A[x]$, 
 then we have $x_{ij}\neq 0, x_{k\ell}\neq 0$ and since $a_{i\ell}\wedge a_{kj}=0$ there are three options for the $2\times 2$ submatrix of $A[x]$ containing these entries
\[
\begin{bmatrix}
    x_{ij}& 0\\
    0 & x_{k\ell}
\end{bmatrix} \text{ or } 
\begin{bmatrix}
    x_{ij}& 0\\
    x_{kj} & x_{k\ell}
\end{bmatrix} \text{ or }
\begin{bmatrix}
    x_{ij}& x_{i\ell}\\
    0 & x_{k\ell}
\end{bmatrix}.
\]
In every case the determinant is $x_{ij}x_{k\ell}$. Hence, $x_{ij}x_{k\ell}\in I_2(A)$, as desired. 

On the other hand suppose $x_{ij}x_{k\ell}\in I_2(A)$. In particular this conveys that $x_{ij},x_{k\ell}\in k[A]$ so that $a_{ij}=a_{k\ell}=1$. 
Suppose towards a  contradiction that $\{a_{ij},a_{k\ell}\}$ is not an isolated pair, 
then the following is  a $2\times 2$ submatrix of $A[x]$\[
\begin{bmatrix}
    x_{ij} & x_{i\ell}\\
    x_{kj}& x_{k\ell}
\end{bmatrix}.
\]
Since $x_{ij}x_{k\ell}\in I_2(A)$, this monomial is a linear combination of determinants of $2\times 2$ submatrices of $A[x]$ including the determinant of the submatrix shown above.
However, the term $x_{kj}x_{i\ell}$ in the determinant of the above matrix appears in no other $2\times 2$ minor of $A[x]$ and so it cannot cancel in the linear combination. This yields a contradiction. Thus it must be that $\{a_{ij},a_{k\ell}\}$ is an isolated pair.
%That is, $x_{ij}x_{k\ell}=x_{ij}x_{k\ell}-x_{kj}x_{i\ell}$. So we must have that $a_{ij}=a_{k\ell}= 1$ and at least one of $a_{kj}, a_{i\ell}$ must be zero. Either way, $a_{kj}a_{i\ell}=0$. Thus, $\{a_{ij},a_{k\ell}\}$ are by definition an isolated pair.
\end{proof}

\begin{theorem}\label{thm: monomials=isolated}
        Let $A$ be an $m\times n$ Boolean matrix. If for some integer $s\geq 2$ the monomial $x_{i_1 j_1}\cdots x_{i_sj_s}\in I_s(A)$, then  $\{a_{i_k j_k}\mid 1\leq k\leq s\}$ is an isolated set.
\end{theorem}

\begin{proof}
Proceed by induction on the degree of the monomial generator. 

\noindent \underline{Base Case:} the case $s=2$ is covered in \Cref{lem: 2x2 minors}.
%Suppose $x_{ij}x_{k\ell}\in \text{gens}(I_2(A[x]))$. Then we know it must be the determinant of the $2\times 2$ submatrix 
%\[
%\begin{bmatrix}
%    x_{i,j} & x_{i,\ell}\\
%    x_{k, j}& x_{k,\ell}
%\end{bmatrix}.
%\]
%That is, $x_{ij}x_{k\ell}=x_{ij}x_{k\ell}-x_{kj}x_{i\ell}$. So we must have that $a_{ij}=a_{k\ell}= 1$ and at least one of $a_{kj}, a_{i\ell}$ must be zero. Either way, $a_{kj}a_{i\ell}=0$. Thus, $\{a_{ij},a_{k\ell}\}$ are by definition an isolated pair.\\

\noindent \underline{Induction Hypothesis:} Assume that if $x_{i_1 j_1}\cdots x_{i_{s-1}j_{s-1}}\in I_{s-1}(A)$, then $\{a_{i_1 j_1}\cdots a_{i_{s-1}j_{s-1}}\}$ is an isolated set. Additionally assume that $s\geq 3$.

Suppose we know that $x_{i_1 j_1}\cdots x_{i_sj_s}\in I_s(A)$. This implies that the given monomial is a $k$-linear combination of $s\times s$ minors of $A[x]$. Since the sets of monomials that appear in the expressions of any two $s\times s$ minors of $A[x]$ are disjoint, it follows that there are no possible cancellations in any linear combination, including within a single determinant. Hence $x_{i_1 j_1}\cdots x_{i_sj_s}\in I_s(A)$ is in fact equal to the minor of $A[x]$ in rows $\{i_1, \ldots, i_s\}$ and columns $\{j_1, \ldots, j_s\}$. Denote this minor by $\text{det}_s$. 
We have established
\[
\text{det}_s = x_{i_1 j_1}\cdots x_{i_sj_s}.
\]
Further denote the minor of $A[x]$ in rows $\{i_1, \ldots, i_{s-1}\}$ and columns $\{j_1, \ldots, j_{s-1}\}$ by $\text{det}_{s-1}$.
In particular, this requires that each $x_{i_k,j_k} \in k[A]$ and hence $a_{i_k,j_k}=1$.

In the following considerations we write $x_{i,j}$ for the entry of $A[x]$ in row $i$ and column $j$ regardless of whether this entry is zero or a variable in $k[A]$. By the definition of determinant, we have 
\begin{align}
\text{det}_s  = \sum_{\sigma \in S_s}\sgn(\sigma)x_{i_1 j_{\sigma(1)}}\cdots x_{i_sj_{\sigma(s)}} \label{eq: 1}\\
\text{det}_{s-1}= \sum_{\sigma \in S_{s-1}}\sgn(\sigma)x_{i_1 j_{\sigma(1)}}\cdots x_{i_{s-1}j_{\sigma(s-1)}},  \label{eq:2}
\end{align}
where $\sgn(\sigma)$ denotes the signature of $\sigma.$
Multiplying \eqref{eq:2} by $x_{i_s j_s}$, we get the following expression
\[
\sum_{\sigma \in S_{s-1}}\sgn(\sigma)x_{i_1 j_{\sigma(1)}}\cdots x_{i_{s-1}j_{\sigma(s-1)}}\cdot x_{i_s j_s},
\]
which is a part of \eqref{eq: 1}. But we assumed that \eqref{eq: 1} is equal to $x_{i_1 j_1}\cdots x_{i_sj_s}$, so we must have that every term that is \textbf{not} $x_{i_1 j_1}\cdots x_{i_sj_s}$ is zero. In particular, $x_{i_1 j_{\sigma(1)}}\cdots x_{i_{s-1}j_{\sigma(s-1)}}\cdot x_{i_s j_s}=0$, unless $\sigma$ is the identity permutation. Since $k[x_{11},x_{12},\ldots, x_{ss}]$ is a domain and $x_{i_s,j_s}\neq 0$ we deduce that $x_{i_1 j_{\sigma(1)}}\cdots x_{i_{s-1}j_{\sigma(s-1)}}=0$, for all $\sigma$ that are not the identity permutation. Therefore  $\text{det}_{s-1}=x_{i_1 j_1}\cdots x_{i_{s-1}j_{s-1}}\in I_{s-1}(A)$.
The induction hypothesis now yields the following isolated set 
$T_{s}=\{a_{i_1,j_1},\ldots, a_{i_{s-1},j_{s-1}}\}$.



%Thus we are left with $x_{i_1 j_1}\cdots x_{i_{s-1}j_(s-1)}$ as the only nonzero term in $(2)$. Hence we have 
%\[
%\text{det}_{s-1}(A[x]) =x_{i_1 j_1}\cdots x_{i_{s-1}j_(s-1)}.
%\]
The same argument, applied for any $1\leq \ell \leq s$, yields that the set
\[
T_\ell=\{a_{i_1 j_1}\cdots \widehat{a_{i_\ell j_\ell}} \cdots a_{i_{s}j_{s}}\}
\]
is an isolated set. Since any pair $\{a_{i_p j_p},a_{i_qj_q}\}$  with $1\leq p \leq s$, $1\leq q\leq s$ and $q\neq p$, is contained in at least one of the sets  $T_\ell$, $\{a_{i_p j_p},a_{i_qj_q}\}$  is an isolated pair for all $q$ and $p$. Therefore we have that $\{a_{i_1 j_1}\cdots a_{i_{s}j_{s}}\}$ is an isolated set, as desired.
\end{proof}

\end{document}
